\section{Síntesis y ampliación}

La contribución central de NLLB puede condensarse en una sola frase: \textbf{la traducción multilingüe a gran escala es un trabajo de ingeniería de sistemas que diseñan en conjunto las necesidades centradas en las personas, la infraestructura de evaluación, la producción de datos, la capacidad del modelo, la planificación del entrenamiento y la validación de seguridad.} Si solo se recuerda MoE o el 44\% de BLEU, se pierde la metodología que el curso realmente quiere transmitir.

\subsection{Doce conclusiones aplicables}

\begingroup
\footnotesize
\setlength{\columnsep}{1.2em}
\begin{multicols}{2}
\begin{enumerate}[itemsep=0.08em,topsep=0.15em,parsep=0pt,leftmargin=*]
  \item Primero confirmar las necesidades de la comunidad; después decidir idiomas, sistemas de escritura, direcciones y escenarios.
  \item Definir high/low-resource como una condición de los datos, no como un orden por población o por valor.
  \item Antes de entrenar, construir un evaluation set comparable, traducido por personas y que abarque varias direcciones.
  \item Reservar el costo real de los estándares lingüísticos, los traductores, los reviewers y el post-editing.
  \item Usar una pequeña cantidad de seed confiable para calibrar LID, el encoder, los filter y la MT inicial.
  \item Diseñar el web mining como un pipeline iterativo y rastreable, no como una extracción única.
  \item Registrar para cada conjunto de datos su provenance, noise, scale y teacher dependency.
  \item El balance de la cola larga debe aumentar la diversidad efectiva; no basta con duplicar las muestras de bajos recursos.
  \item La ganancia de capacidad de MoE debe evaluarse junto con el sobreajuste en bajos recursos, el enrutamiento y la regularización.
  \item Reportar los resultados por direction, resource group y metric, no solo el promedio.
  \item La evaluación automatic, la human y la de toxicity responden preguntas distintas.
  \item Solo al abrir el benchmark, las herramientas de datos y los modelos puede la comunidad seguir corrigiendo las brechas de cobertura.
\end{enumerate}
\end{multicols}
\endgroup

\subsection{Panorama final del sistema}

Al unir toda la clase se obtiene un ciclo de ingeniería con retroalimentación. Este ciclo no es un proceso lineal en cascada, porque la evaluación humana, las fallas de seguridad y la retroalimentación de la comunidad obligan al equipo a regresar a las etapas de estándares, datos y modelo:
\[
\begin{aligned}
\text{Community needs}
&\rightarrow\text{Standards \& Evaluation}
\rightarrow\text{Seed Data}
\rightarrow\text{Mining / BT}\\
&\rightarrow\text{MoE Training}
\rightarrow\text{Human \& Safety Validation}
\rightarrow\text{Community feedback}.
\end{aligned}
\]
Cada flecha puede introducir un sesgo de selección y cada retroalimentación puede modificar los estándares de las etapas previas. Un sistema verdaderamente maduro no se completa con un solo entrenamiento: es capaz de explicar de dónde provienen los datos, por qué el modelo falla en ciertos idiomas, qué métricas respaldan las conclusiones y quién tiene la facultad de corregir los errores.

\begin{importantbox}{Criterio de juicio final}
«No Language Left Behind» no puede declararse terminado al alcanzar cierto número de idiomas. Un juicio más confiable consiste en preguntar: ¿los idiomas de la cola cuentan con datos y evaluaciones auditables?, ¿los hablantes nativos participan en las necesidades, los estándares y la validación?, ¿el sistema reporta la calidad por dirección y los errores de alto riesgo?, ¿los recursos abiertos permiten que la comunidad reproduzca, critique y mejore? La cifra de cobertura solo puede ser una de las piezas de esta evidencia.
\end{importantbox}

\subsection{Lecturas adicionales}

\begingroup
\scriptsize
\begin{itemize}[itemsep=0pt,topsep=0.15em,parsep=0pt]
  \item NLLB Team, \href{https://arxiv.org/abs/2207.04672}{No Language Left Behind: Scaling Human-Centered Machine Translation}: informe completo sobre datos, modelos, evaluación e impacto social.
  \item Meta AI, \href{https://research.facebook.com/publications/no-language-left-behind/}{NLLB publication page} y \href{https://github.com/facebookresearch/fairseq/tree/nllb}{fairseq NLLB release}: punto de entrada oficial al proyecto y al código de los modelos.
  \item \href{https://github.com/facebookresearch/flores}{FLORES}: datos de evaluación multilingüe; \href{https://github.com/facebookresearch/stopes}{Stopes}: herramienta escalable de minería de datos multilingües.
  \item Schwenk et al., \href{https://arxiv.org/abs/1907.05791}{WikiMatrix} y \href{https://github.com/facebookresearch/LASER}{LASER}: fundamentos de la minería de corpus mediante vectores de oraciones.
\end{itemize}
\endgroup

\end{document}
